\documentclass[11pt,a4paper]{article}

\usepackage[margin=2.5cm, headheight=14pt]{geometry}
\usepackage{booktabs}
\usepackage{multirow}
\usepackage{array}
\usepackage{xcolor}
\usepackage{hyperref}
\usepackage{amsmath}
\usepackage{microtype}
\usepackage{parskip}
\usepackage{enumitem}
\usepackage{caption}
\usepackage{fancyhdr}
\usepackage{titlesec}
\usepackage{tabularx}
\usepackage{colortbl}

\definecolor{ownerblue}{RGB}{30,80,160}
\definecolor{bertcol}{RGB}{210,230,255}
\definecolor{llmcol}{RGB}{210,255,220}
\definecolor{s2tollcol}{RGB}{255,215,215}
\definecolor{s2tmmrcol}{RGB}{255,235,200}
\definecolor{bestcol}{RGB}{255,255,180}

\hypersetup{colorlinks=true,linkcolor=ownerblue,urlcolor=ownerblue,citecolor=ownerblue}

\pagestyle{fancy}
\fancyhf{}
\rhead{\textcolor{gray}{\small OWNER \& Span2Type -- Cluster Naming Evaluation}}
\lhead{\textcolor{gray}{\small BERTScore Report (Version~6)}}
\rfoot{\thepage}

\titleformat{\section}{\large\bfseries\color{ownerblue}}{{\thesection}}{1em}{}
\titleformat{\subsection}{\normalsize\bfseries}{{\thesubsection}}{1em}{}

\begin{document}

\begin{center}
  {\LARGE\bfseries\color{ownerblue}
    Cluster Naming Evaluation Report \normalsize\textit{(Version~6)}\\[0.4em]
    \large OWNER vs.\ Span2Type: BERTScore Comparison}\\[0.8em]
  {\large Five CrossNER Domains}\\[0.6em]
\end{center}
\vspace{0.4em}\hrule\vspace{1.2em}

\begin{abstract}
This report compares four cluster-naming methods across five CrossNER domains
(AI, Literature, Music, Politics, Science) using BERTScore (RoBERTa-large)
against confusion-matrix-derived ground-truth labels,
All scores are measured at the \textbf{Entity Typing (ET) stage}.
Two methods belong to \textbf{OWNER}: BERT-based MLM voting (avg
F$_1{=}0.897$) and LLM-based naming via Ollama (avg F$_1{=}0.908$).
Two methods belong to \textbf{Span2Type}, both using Maximum Marginal
Relevance (MMR) exemplar selection:
MMR-LLM naming (avg F$_1{=}0.924$, best overall) and MMR-BERT naming
(avg F$_1{=}0.896$).
All scores use the corrected ground truth established in Version~5.
S2T-MMR-LLM is the strongest method; LLM-based methods consistently
outperform MLM/BERT voting across all domains.
\end{abstract}

%─── 1. Introduction ──────────────────────────────────────────────────────────
\section{Introduction}

Unsupervised NER systems such as OWNER~\cite{genest2025owner} group entity
spans into clusters labelled by integer IDs.
Cluster naming assigns human-readable type labels without manual annotation.
This report evaluates four naming strategies:

\begin{itemize}[leftmargin=1.5em]
  \item \textbf{OWNER-BERT}: for each entity, fill
    \texttt{``\{sentence\} \{entity\} is a [MASK].''} with BERT-base-uncased;
    the most frequent top-1 token across the cluster is the name.
  \item \textbf{OWNER-LLM}: sample $n{=}16$ entities per cluster (uniform),
    query Llama-3.1-8B via Ollama (Figure~3 prompt, temperature~$=0$).
  \item \textbf{S2T-MMR-LLM}: uses \textit{Maximum Marginal Relevance}
    (MMR) to select $k{=}5$ diverse, centroid-representative entities per
    cluster, then queries Llama-3.1-8B via Ollama with full example sentences
    (richer context); produces multi-word descriptive names.
  \item \textbf{S2T-MMR-BERT}: same MMR exemplar selection ($k{=}5$) as
    S2T-MMR-LLM, but uses BERT-base-uncased MLM filling instead of an LLM
    — the top-1 token from \texttt{``\{sentence\} \{entity\} is a [MASK].''}
    is the cluster name.
\end{itemize}

\noindent
All methods are evaluated against ground-truth labels derived from the
column-wise argmax of the entity-type confusion matrix
($\operatorname{argmax}_i\,C_{ij}$), consistent with the corrected
methodology of Version~5.

%─── 2. Experimental Setup ────────────────────────────────────────────────────
\section{Experimental Setup}

\subsection{MLflow Run IDs}

\begin{table}[h]
\centering
\caption{MLflow run IDs for OWNER and Span2Type experiments.}
\label{tab:runs}
\setlength{\tabcolsep}{4pt}
\small
\begin{tabular}{l l l l l}
\toprule
\textbf{Domain}
  & \textbf{OWNER}
  & \textbf{S2T-MMR-LLM}
  & \textbf{S2T-MMR-BERT}
  & \textbf{S2T-no-MMR} \\
\midrule
AI
  & \texttt{ea6e40f8...eb6e}
  & \texttt{6b105207...2b2}
  & \texttt{9e28895d...ac9}
  & \texttt{22d48750...771} \\
Literature
  & \texttt{fb1215f7...f0}
  & \texttt{29e11a34...be}
  & \texttt{89ec5ac6...68b}
  & \texttt{a24b0559...dd} \\
Music
  & \texttt{2c886b9b...c5}
  & \texttt{56ad65e3...dc}
  & \texttt{d70574bc...b2}
  & \texttt{08d4cdde...2f} \\
Politics
  & \texttt{9c02f083...6d}
  & \texttt{610c691d...7b}
  & \texttt{5b3e1bf3...12}
  & \texttt{56aacc85...40} \\
Science
  & \texttt{f72db65d...96}
  & \texttt{efcfb115...55}
  & \texttt{eec16336...fe}
  & \texttt{14b05751...2b} \\
\bottomrule
\end{tabular}
\end{table}

\subsection{Ground-Truth Alignment}

For OWNER runs, GT labels are read from \texttt{et\_test\_cluster\_naming\_gt.json}
(confusion-matrix argmax, corrected in Version~5).
For Span2Type runs, GT labels are derived at runtime from
\texttt{et\_test\_confusion-matrix.pt} using the same argmax procedure.

\subsection{Cluster Statistics}

\begin{table}[h]
\centering
\caption{Number of clusters per domain and pipeline stage.}
\label{tab:clusters}
\setlength{\tabcolsep}{6pt}
\small
\begin{tabular}{l r r r}
\toprule
\textbf{Domain} & \textbf{True $k$} & \textbf{OWNER / S2T-v1} & \textbf{S2T-MMR-BERT (NER)} \\
\midrule
AI         & 14 & 16 & 14 \\
Literature & 12 & 14 & 12 \\
Music      & 13 & 20 & 20 \\
Politics   &  9 & 22 & 22 \\
Science    & 17 & 20 & 20 \\
\midrule
Total      & 65 & 92 & 88 \\
\bottomrule
\end{tabular}
\end{table}

%─── 3. Results ───────────────────────────────────────────────────────────────
\section{Results}

\subsection{BERTScore F\textsubscript{1} Comparison}

Table~\ref{tab:f1} presents BERTScore F$_1$ for all four methods
at the \textbf{Entity Typing (ET) stage}.
The best score per domain is shown in \textbf{bold}.

\begin{table}[h]
\centering
\caption{BERTScore F$_1$ at the ET stage -- four cluster-naming methods
across five CrossNER domains. Best per domain in
\textbf{bold}.}
\label{tab:f1}
\setlength{\tabcolsep}{5pt}
\small
\begin{tabular}{l cccc}
\toprule
\textbf{Domain}
  & \cellcolor{bertcol}\textbf{OW-BERT}
  & \cellcolor{llmcol}\textbf{OW-LLM}
  & \cellcolor{s2tollcol}\textbf{S2T-MMR-LLM}
  & \cellcolor{s2tmmrcol}\textbf{S2T-MMR-BERT} \\
\midrule
AI         & 0.900          & 0.892          & \textbf{0.922} & 0.900 \\
Literature & 0.891          & \textbf{0.923} & 0.913          & 0.889 \\
Music      & 0.906          & 0.928          & \textbf{0.943} & 0.911 \\
Politics   & 0.892          & 0.910          & \textbf{0.928} & 0.887 \\
Science    & 0.894          & 0.886          & \textbf{0.912} & 0.893 \\
\midrule
\textbf{Avg} & 0.897        & 0.908          & \textbf{0.924} & 0.896 \\
\bottomrule
\end{tabular}
\end{table}

\noindent
Overall ranking (ET stage):
\begin{center}
S2T-MMR-LLM ($\mathbf{0.924}$)
$>$ OW-LLM ($0.908$)
$>$ OW-BERT $\approx$ S2T-MMR-BERT ($0.897\,/\,0.896$)
\end{center}

\subsection{S2T-MMR-BERT Detailed Scores (P\,/\,R\,/\,F\textsubscript{1})}

Table~\ref{tab:mmr_detail} shows full BERTScore Precision, Recall, and
F$_1$ for S2T-MMR-BERT at the Entity Typing (ET) stage.

\begin{table}[h]
\centering
\caption{S2T-MMR-BERT BERTScore at the ET stage (P\,/\,R\,/\,F$_1$).}
\label{tab:mmr_detail}
\setlength{\tabcolsep}{6pt}
\small
\begin{tabular}{l ccc}
\toprule
& \multicolumn{3}{c}{\cellcolor{s2tmmrcol}\textbf{ET Stage}} \\
\cmidrule(lr){2-4}
\textbf{Domain} & P & R & F$_1$ \\
\midrule
AI         & 0.897 & 0.904 & 0.900 \\
Literature & 0.883 & 0.896 & 0.889 \\
Music      & 0.906 & 0.916 & 0.911 \\
Politics   & 0.880 & 0.894 & 0.887 \\
Science    & 0.893 & 0.893 & 0.893 \\
\midrule
\textbf{Avg} & 0.892 & 0.900 & 0.896 \\
\bottomrule
\end{tabular}
\end{table}

\subsection{S2T-MMR-LLM Detailed Scores (P\,/\,R\,/\,F\textsubscript{1})}

Table~\ref{tab:mmr_llm_detail} shows full BERTScore Precision, Recall, and
F$_1$ for S2T-MMR-LLM at the Entity Typing (ET) stage.

\begin{table}[h]
\centering
\caption{S2T-MMR-LLM BERTScore at the ET stage (P\,/\,R\,/\,F$_1$).}
\label{tab:mmr_llm_detail}
\setlength{\tabcolsep}{6pt}
\small
\begin{tabular}{l ccc}
\toprule
& \multicolumn{3}{c}{\cellcolor{s2tollcol}\textbf{ET Stage}} \\
\cmidrule(lr){2-4}
\textbf{Domain} & P & R & F$_1$ \\
\midrule
AI         & 0.919 & 0.925 & 0.922 \\
Literature & 0.909 & 0.917 & 0.913 \\
Music      & 0.940 & 0.946 & 0.943 \\
Politics   & 0.925 & 0.931 & 0.928 \\
Science    & 0.913 & 0.911 & 0.912 \\
\midrule
\textbf{Avg} & 0.921 & 0.926 & 0.924 \\
\bottomrule
\end{tabular}
\end{table}

\subsection{Qualitative Naming Examples}

Table~\ref{tab:examples} shows representative S2T-MMR-BERT predictions
with per-cluster BERTScore F$_1$.

\begin{table}[h]
\centering
\caption{Selected S2T-MMR-BERT cluster naming examples.
\textcolor{green!60!black}{Green} = exact/near-exact match;
\textcolor{orange!70!black}{orange} = semantically close;
\textcolor{red!70!black}{red} = incorrect or noise.}
\label{tab:examples}
\setlength{\tabcolsep}{4pt}
\small
\begin{tabularx}{\textwidth}{l l l l r}
\toprule
\textbf{Dom.} & \textbf{GT} & \textbf{Prediction} & \textbf{Quality} & \textbf{F$_1$} \\
\midrule
AI   & university    & \textcolor{green!60!black}{university}   & Exact   & 1.000 \\
AI   & conference    & \textcolor{green!60!black}{conference}   & Exact   & 1.000 \\
AI   & task          & \textcolor{orange!70!black}{problem}     & Close   & 0.943 \\
AI   & algorithm     & \textcolor{orange!70!black}{standard}    & Partial & 0.908 \\
AI   & researcher    & \textcolor{orange!70!black}{physicist}   & Close   & 0.858 \\
AI   & programlang   & \textcolor{red!70!black}{standard}       & Off     & 0.838 \\
AI   & person        & \textcolor{red!70!black}{vegetarian}     & Noise   & 0.853 \\
\midrule
Lit. & book          & \textcolor{orange!70!black}{novel}       & Close   & 0.906 \\
Lit. & event         & \textcolor{orange!70!black}{festival}    & Close   & 0.918 \\
Lit. & country       & \textcolor{orange!70!black}{nation}      & Close   & 0.924 \\
Lit. & writer        & \textcolor{orange!70!black}{poet}        & Close   & 0.897 \\
Lit. & award         & \textcolor{orange!70!black}{retrospective} & Partial & 0.865 \\
Lit. & book          & \textcolor{red!70!black}{vegetarian}     & Noise   & 0.843 \\
\midrule
Mus. & country       & \textcolor{green!60!black}{country}      & Exact   & 1.000 \\
Mus. & band          & \textcolor{green!60!black}{band}         & Exact   & 1.000 \\
Mus. & musicgenre    & \textcolor{orange!70!black}{genre}       & Close   & 0.936 \\
Mus. & musicalartist & \textcolor{orange!70!black}{musician}    & Close   & 0.930 \\
Mus. & award         & \textcolor{orange!70!black}{nomination}  & Close   & 0.897 \\
Mus. & band          & \textcolor{red!70!black}{vegetarian}     & Noise   & 0.848 \\
\midrule
Pol. & politicalparty & \textcolor{orange!70!black}{party}      & Close   & 0.930 \\
Pol. & location      & \textcolor{orange!70!black}{city}        & Close   & 0.937 \\
Pol. & politician    & \textcolor{orange!70!black}{democrat}    & Close   & 0.899 \\
Pol. & person        & \textcolor{red!70!black}{vegetarian}     & Noise   & 0.853 \\
\midrule
Sci. & university    & \textcolor{green!60!black}{university}   & Exact   & 1.000 \\
Sci. & protein       & \textcolor{green!60!black}{protein}      & Exact   & 1.000 \\
Sci. & scientist     & \textcolor{orange!70!black}{physicist}   & Close   & 0.935 \\
Sci. & chemicalcompound & \textcolor{orange!70!black}{catalyst} & Partial & 0.881 \\
Sci. & chemicalelement & \textcolor{orange!70!black}{semiconductor} & Partial & 0.834 \\
Sci. & person        & \textcolor{red!70!black}{vegetarian}     & Noise   & 0.853 \\
\bottomrule
\end{tabularx}
\end{table}

%─── 4. MMR Ablation ──────────────────────────────────────────────────────────
\section{MMR Ablation}

To isolate the effect of Maximum Marginal Relevance exemplar selection,
we compare two Span2Type configurations:
\textbf{no-MMR} ($k{=}16$) versus
\textbf{MMR} ($k{=}5$ diverse exemplars).
The no-MMR run IDs are listed in the rightmost column of
Table~\ref{tab:runs}.

\subsection{Ollama Backend}

\begin{table}[h]
\centering
\caption{S2T Ollama: no-MMR ($k{=}16$, no lem) vs MMR ($k{=}5$, +lem).
Best per domain in \textbf{bold}.}
\label{tab:ablation_oll}
\setlength{\tabcolsep}{6pt}
\small
\begin{tabular}{l cc r}
\toprule
\textbf{Domain}
  & \cellcolor{s2tollcol}\textbf{no-MMR}
  & \cellcolor{s2tmmrcol}\textbf{MMR}
  & \textbf{$\Delta$} \\
\midrule
AI         & 0.9218          & 0.9219          & $+0.0001$ \\
Literature & 0.9063          & \textbf{0.9197} & $+0.0134$ \\
Music      & \textbf{0.9539} & 0.9509          & $-0.0030$ \\
Politics   & \textbf{0.9302} & 0.9293          & $-0.0009$ \\
Science    & 0.9084          & \textbf{0.9156} & $+0.0072$ \\
\midrule
\textbf{Avg} & 0.9241        & \textbf{0.9275} & $+0.0034$ \\
\bottomrule
\end{tabular}
\end{table}

\noindent
With the Ollama backend, MMR improves average F$_1$ by
$+0.003$ overall, but the gain is domain-dependent.
Literature and Science benefit most ($+0.013$, $+0.007$), where the
no-MMR uniform sample yields redundant entities that lead the LLM to
produce vague labels (\textit{Book}, \textit{Science}); MMR selects
diverse centroid-representative exemplars, eliciting more specific
outputs (\textit{Literary work}, \textit{Scientific discipline}).
Conversely, Music and Politics show a small regression ($-0.003$,
$-0.001$): no-MMR already samples a broad mix of band names, song
titles, and politician names that give the LLM sufficient signal,
while MMR occasionally over-concentrates on centroid entities,
slightly narrowing the context.
The AI domain is effectively tied ($+0.0001$), consistent with
abstract technical types being hard for both sampling strategies.

\subsection{BERT Backend}

\begin{table}[h]
\centering
\caption{S2T BERT: no-MMR ($k{=}16$) vs MMR ($k{=}5$).
; $\Delta$ reflects MMR.
Best per domain in \textbf{bold}.}
\label{tab:ablation_bert}
\setlength{\tabcolsep}{6pt}
\small
\begin{tabular}{l cc r}
\toprule
\textbf{Domain}
  & \cellcolor{bertcol}\textbf{no-MMR}
  & \cellcolor{s2tmmrcol}\textbf{MMR}
  & \textbf{$\Delta$} \\
\midrule
AI         & \textbf{0.9037} & 0.9000          & $-0.0037$ \\
Literature & 0.8874          & \textbf{0.8888} & $+0.0014$ \\
Music      & 0.8948          & \textbf{0.9110} & $+0.0162$ \\
Politics   & 0.8830          & \textbf{0.8868} & $+0.0038$ \\
Science    & \textbf{0.8935} & 0.8925          & $-0.0010$ \\
\midrule
\textbf{Avg} & 0.8925        & \textbf{0.8958} & $+0.0033$ \\
\bottomrule
\end{tabular}
\end{table}

\noindent
With the BERT backend, MMR yields a small but consistent improvement on average ($+0.003$),
with the largest gain in Music ($+0.016$): the no-MMR uniform sample
includes many redundant entity mentions from a few dominant clusters
(e.g., several occurrences of the same band name), causing the MLM
head to converge on generic tokens; MMR diversifies the exemplar set,
shifting the argmax towards more cluster-specific tokens
(\textit{band}, \textit{musician}, \textit{genre}).
AI and Science show minor regressions ($-0.004$, $-0.001$) because
MMR's centroid-representative selection occasionally surfaces
technically-worded entity spans where the MLM's most frequent token
is less aligned with the GT label than the token obtained under
uniform sampling.
Crucially, MMR cannot suppress the \texttt{vegetarian} artefact:
the noise prediction persists in 3--5 clusters per domain regardless
of exemplar selection strategy, confirming that the limitation is
intrinsic to greedy MLM top-1 decoding.

\subsection{Ablation Summary}

\begin{table}[h]
\centering
\caption{MMR ablation summary: lowest (no-MMR) vs highest (MMR) per backend.}
\label{tab:ablation_summary}
\setlength{\tabcolsep}{8pt}
\small
\begin{tabular}{l cc r}
\toprule
\textbf{Backend} & \textbf{no-MMR} & \textbf{MMR} & \textbf{$\Delta$} \\
\midrule
Ollama & 0.9241 & \textbf{0.9275} & $+0.0034$ \\
BERT   & 0.8925 & \textbf{0.8958} & $+0.0033$ \\
\bottomrule
\end{tabular}
\end{table}

Both backends improve by ${\approx}+0.003$ from lowest to highest configuration.
For the Ollama backend the gain requires MMR selection, MMR elicits richer, more plural-inflected outputs from
the LLM; .
For the BERT backend, the $+0.003$ gain comes entirely from MMR selection concentrating the MLM
prompt on diverse, centroid-representative exemplars.

%─── 5. Analysis ──────────────────────────────────────────────────────────────
\section{Analysis}

\subsection{S2T-MMR-LLM Is the Strongest Method}

S2T-MMR-LLM (avg F$_1{=}0.924$) outperforms all other methods in four
of five domains. It combines MMR exemplar selection with rich sentence context
passed to Llama-3.1-8B, enabling the LLM to infer fine-grained types such as
\textit{machine learning library} (GT: product), \textit{legislative body}
(GT: organisation), and \textit{biological processes} (GT: misc).

\subsection{LLM Back-end Drives the Quality Gap, Not MMR Alone}

Both S2T methods use the same MMR exemplar selection ($k{=}5$), yet
S2T-MMR-LLM (0.924) substantially outperforms S2T-MMR-BERT (0.896).
This $+0.028$ gap is attributable entirely to the naming back-end:
LLM outputs are multi-word, context-aware descriptions, whereas BERT
top-1 token voting collapses to a single common token.
Comparing S2T-MMR-BERT with OWNER-BERT (0.897) — the latter without MMR —
shows that MMR selection itself gives no measurable benefit ($-0.001$),
confirming that the bottleneck for MLM naming lies in the vocabulary
and greedy decoding of the MLM head, not in exemplar selection.

\subsection{LLM Methods Outperform MLM Voting}

Both LLM methods (OW-LLM: 0.908, S2T-MMR-LLM: 0.924) score above both
MLM methods (OW-BERT: 0.897, S2T-MMR-BERT: 0.896).
MLM top-1 voting is limited to single-token outputs and frequently falls
back to degenerate predictions (\texttt{vegetarian}) for ambiguous clusters.
LLM methods produce multi-word names with stronger semantic alignment
to CrossNER entity type labels.

\subsection{The \texttt{vegetarian} Artefact}

The \texttt{vegetarian} noise token appears consistently across all MLM
methods: 4--5 clusters per domain (out of 14--22) receive this prediction.
It is the most impactful single failure mode for BERT/MLM-based naming,
producing per-cluster F$_1 \approx 0.85$ against GT labels such as
\texttt{person}, \texttt{book}, or \texttt{band}.
LLM methods are immune to this artefact.

\subsection{Per-Domain Observations}

\begin{itemize}[leftmargin=1.5em]
  \item \textbf{AI}: All methods score 0.89--0.92. Abstract types
    (\texttt{algorithm}, \texttt{programlang}, \texttt{task}) are hard for
    all methods; only S2T-MMR-LLM captures \textit{programming languages}
    and \textit{machine learning} as reasonable proxies.
  \item \textbf{Literature}: OW-LLM wins (0.923) due to descriptive names
    (\textit{Novel}, \textit{Author}, \textit{Film Festival}) that align
    well with CrossNER types.
  \item \textbf{Music}: S2T-MMR-LLM dominates (0.943). Three exact matches
    for S2T-MMR-BERT (\texttt{band}, \texttt{band}, \texttt{country})
    keep it competitive (0.911).
  \item \textbf{Politics}: S2T-MMR-LLM best (0.928) via \textit{political
    party}, \textit{legislative body}, \textit{politician}. BERT methods
    predict generic tokens (\textit{party}, \textit{democrat}).
  \item \textbf{Science}: S2T-MMR-LLM leads (0.912); S2T-MMR-BERT recovers
    with exact matches (\texttt{university}, \texttt{protein}).
    Science-specific compound types (\texttt{chemicalelement},
    \texttt{academicjournal}) remain hard for all methods.
\end{itemize}

%─── 6. Conclusion ────────────────────────────────────────────────────────────
\section{Conclusion}

The four-method ranking at the ET stage is:

\begin{center}
S2T-MMR-LLM ($\mathbf{0.924}$)
$>$ OW-LLM ($0.908$)
$>$ OW-BERT ($0.897$)
$\approx$ S2T-MMR-BERT ($0.896$)
\end{center}

\noindent
The MMR ablation (Section~4) shows that both backends gain
${\approx}{+}0.003$ when moving from the no-MMR baseline to the best
MMR.
Across all conditions, the primary quality driver remains the naming
back-end: LLM methods outperform MLM methods by $0.028$--$0.032$
regardless of exemplar selection strategy.
The \texttt{vegetarian} artefact remains the main failure mode for all
MLM/BERT approaches and cannot be resolved without filtering heuristics
or switching to an LLM-based back-end.

\begin{thebibliography}{9}

\bibitem{genest2025owner}
P.-Y.~Genest, P.-E.~Portier, E.~Egyed-Zsigmond, and M.~Lovisetto,
``OWNER -- Toward Unsupervised Open-World Named Entity Recognition,''
\textit{IEEE Access}, vol.~13, pp.~50077--50096, 2025.

\bibitem{zhang2020bertscore}
T.~Zhang, V.~Kishore, F.~Wu, K.~Q.~Weinberger, and Y.~Artzi,
``BERTScore: Evaluating Text Generation with BERT,''
in \textit{Proc. ICLR}, 2020.

\end{thebibliography}

\end{document}
